\documentclass[../../../include/open-logic-section]{subfiles}

\begin{document}

\olfileid[pt]{sth}{spine}{valpha}

\olsection{Definindo os estágios como os $V_\alpha$}

No \olref[sth][ordinals][]{chap}, definimos boas ordenações e os
ordinais (de von Neumann). Neste capítulo, usaremos isto para
caracterizar a própria hierarquia de conjuntos. Para tal, recordemos
que na \olref[sth][ordinals][opps]{sec} definimos a noção de
ordinal sucessor e ordinal limite. Usamos estas ideias na definição seguinte:

\begin{defn}\ollabel{defValphas}
	\begin{align*}
	V_\emptyset &\defis \emptyset\\
	V_{\ordsucc{\alpha}} &\defis \Pow{V_\alpha} & & 
	\text{para qualquer ordinal }\alpha\\
	V_{\alpha} &\defis \bigcup_{\gamma < \alpha} V_\gamma & & 
	\text{quando }\alpha\text{ é um ordinal limite}
\end{align*}
\end{defn}
\noindent
Será uma definição por \emph{recursão transfinita} nos
ordinais. A este respeito, convém comparar com definições recursivas
de funções nos números naturais.\footnote{Cf.\ as
definições de adição, multiplicação e exponenciação em
\olref[sfr][infinite][dedekind]{sec}.} Como ao tratar dos números
naturais, define-se um caso base e casos sucessor; mas ao tratar
de ordinais, também precisamos de descrever o comportamento nos casos
\emph{limite}. 

Esta definição dos $V_\alpha$ será um marco importante. Motivamos
informalmente a nossa hierarquia de conjuntos como formando conjuntos por
\emph{estágios}. Os $V_\alpha$ são, em efeito, precisamente esses estágios.
Importa, contudo, que se trata de uma caracterização \emph{interna} dos
estágios. Em vez de sugerir um possível \emph{modelo} da
teoria, teremos definido os estágios \emph{dentro} da nossa teoria dos conjuntos.

\end{document}
